%! Setting Up the SVD — Unit 7, Lesson 7.0 — Homework
\documentclass[10pt]{article}
\usepackage{linalg-article}
\usepackage{linalg-boxes}
\newcommand{\TallMath}[1]{$\displaystyle #1\rule[-1.4em]{0pt}{3.2em}$}
\newcommand{\uu}{\mathbf{u}}
\newcommand{\vv}{\mathbf{v}}
\newcommand{\ee}{\mathbf{e}}
\newcommand{\zero}{\mathbf{0}}
\newcommand{\T}{^{\top}}

\begin{document}

\pageheader{Unit 7, Lesson 7.0}{Homework}
\namedateperiod

\vspace{0.12in}

\begin{notesbox}{Practice}
A matrix stretches the unit circle into an \textbf{ellipse}. The \textbf{singular values}
$\sigma_1\ge\sigma_2\ge0$ are its half-widths; a perpendicular input pair $\vv_1,\vv_2$ maps to a
perpendicular output pair with $A\vv_i=\sigma_i\uu_i$.

\begin{enumerate}[leftmargin=1.9em, itemsep=12pt]
  \item \textbf{Read $\sigma$ off a diagonal matrix.} Give the singular values $\sigma_1,\sigma_2$
        (larger first) and describe the ellipse the unit circle becomes.
    \begin{enumerate}[label=(\alph*), leftmargin=1.8em, itemsep=4pt]
      \item $\begin{bmatrix} 5 & 0 \\ 0 & 2 \end{bmatrix}$ \hfill
      \item $\begin{bmatrix} 1 & 0 \\ 0 & 6 \end{bmatrix}$ \hfill
      \item $\begin{bmatrix} 3 & 0 \\ 0 & 3 \end{bmatrix}$
    \end{enumerate}
    \vspace{1.1cm}
  \item \textbf{Perpendicular in, perpendicular out.} For $S=\begin{bmatrix} 0 & 4 \\ 3 & 0 \end{bmatrix}$,
        compute $S\ee_1$ and $S\ee_2$. Show the outputs are perpendicular, give $\sigma_1,\sigma_2$, and
        state the output direction $\uu_1$ belonging to $\sigma_1$. Is $\uu_1$ the same as its input
        direction? \writelines{2}
  \item \textbf{A symmetric matrix (recall Unit~6).} For $A=\begin{bmatrix} 3 & 1 \\ 1 & 3 \end{bmatrix}$,
        the vectors $\begin{bmatrix} 1 \\ 1 \end{bmatrix}$ and $\begin{bmatrix} 1 \\ -1 \end{bmatrix}$ are
        perpendicular eigenvectors. Compute $A\begin{bmatrix} 1 \\ 1 \end{bmatrix}$ and
        $A\begin{bmatrix} 1 \\ -1 \end{bmatrix}$, and read off the singular values. \writelines{2}
  \item \textbf{Why never negative?} Explain why a singular value can never be negative, and why the
        singular values are listed from largest to smallest. \writelines{2}
  \item \textbf{Why the SVD wins.} In a sentence or two, describe what $A=U\Sigma V\T$ says a matrix does
        (rotate, stretch, rotate), and state one reason the SVD is more widely usable than eigenvectors.
        \writelines{2}
\end{enumerate}
\end{notesbox}

\begin{extensionbox}
\textbf{Finding singular values with $A\T A$ (a first look at Lesson~7.1).} The singular values are
$\sigma=\sqrt{\lambda}$, where $\lambda$ are the eigenvalues of the symmetric matrix $A\T A$. Use
$S=\begin{bmatrix} 0 & 4 \\ 3 & 0 \end{bmatrix}$ from Problem~2.
\begin{enumerate}[label=(\alph*), leftmargin=1.8em, itemsep=5pt]
  \item Compute $A\T A$ (remember $A\T$ swaps rows and columns).
  \item Find the eigenvalues $\lambda$ of $A\T A$, then take $\sigma=\sqrt{\lambda}$. Do they match the
        $\sigma$'s you read off in Problem~2?
\end{enumerate}
\writelines{3}
\end{extensionbox}

\begin{spiralbox}
\textbf{Coming next --- Lesson 7.1: Singular Values and Singular Vectors.} Today you \emph{read} singular
values off nice matrices. Next we \emph{find} them for any matrix: the eigenvalues of $A\T A$ give
$\sigma^2$, and its eigenvectors give the input directions $\vv_i$ --- then $\uu_i=A\vv_i/\sigma_i$.
\end{spiralbox}

\end{document}
